\section{Theoretical analysis}

\newtheorem{adaptationlemma}[theorem]{Lemma candidate}
\newtheorem{adaptationproposition}[theorem]{Proposition candidate}

\subsection{Minimal vector linear specialization}

The following results characterize the immediate and adaptation effects of
a single teacher-induced update. They have explicit algebraic proofs, but
remain preliminary with respect to the complete delayed-feedback system and
do not constitute novelty claims or sequential optimality guarantees.

Under squared loss, consider
\begin{align}
    Y &= (w^*)^\top X + \epsilon,
    &
    F_\theta(x) &= \theta^\top x,
    \\
    D(x) &= (w^*+b)^\top x + \nu,
    &
    e &= \theta-w^*,
    \\
    M &= \mathbb E[XX^\top] \succ 0,
    &
    \alpha &= x^\top e,
    \qquad
    \beta = x^\top b .
\end{align}
The matrix $M$ is the second-moment matrix of the future input distribution
and need not be a centered covariance matrix.

The teacher noise satisfies, conditionally on the current input,
\begin{equation}
    \mathbb E[\nu\mid x]=0,
    \qquad
    \mathbb E[\nu^2\mid x]=\sigma_D^2.
\end{equation}
For the current task-loss comparison, we also assume
\begin{equation}
    \mathbb E[\epsilon\mid x]=0,
    \qquad
    \mathbb E[\epsilon\nu\mid x]=0,
\end{equation}
with finite second moments. These are assumptions of the present
specialization rather than consequences of marginal noise moments. No
Gaussian distribution or full independence assumption is required.

Holding the parameters and the pre-query state fixed, the expected
immediate gain from replacing the cheap response by the teacher response is
\begin{equation}
    \Delta_{\mathrm{now}}(x)
    = \alpha^2-\beta^2-\sigma_D^2,
    \label{eq:immediate-gain}
\end{equation}
where the common irreducible task-noise contribution cancels.

The population excess-risk component of the cheap predictor is
\begin{equation}
    R(e)=e^\top M e.
    \label{eq:population-risk}
\end{equation}
This quantity is not the total system objective: it excludes routing costs,
future operational decisions, and irreducible task noise.

When the teacher is queried, its same operational response is used as an
immediate pseudo-target:
\begin{align}
    \theta^+
    &= \theta+\eta x\bigl(D(x)-\theta^\top x\bigr),
    \\
    e^+
    &= e-\eta x(\alpha-\beta)+\eta x\nu,
    \qquad \eta>0.
    \label{eq:error-update}
\end{align}

\begin{adaptationlemma}[Exact single-update adaptation value]
Under the conditional teacher-noise assumptions above,
\begin{align}
    \Delta_R(x)
    &:= R(e)-\mathbb E[R(e^+)\mid x]
    \notag\\
    &= 2\eta(\alpha-\beta)x^\top M e
    -\eta^2
    \bigl[(\alpha-\beta)^2+\sigma_D^2\bigr]
    x^\top Mx .
    \label{eq:exact-adaptation}
\end{align}
\end{adaptationlemma}

\begin{proof}
Let $d=\alpha-\beta$. From~\eqref{eq:error-update},
\begin{equation}
    e^+=e+\eta x(\nu-d).
\end{equation}
Expanding the quadratic risk gives
\begin{equation}
    R(e^+)
    =
    R(e)
    +2\eta(\nu-d)x^\top Me
    +\eta^2(\nu-d)^2x^\top Mx.
\end{equation}
Conditioning on $x$,
\begin{equation}
    \mathbb E[\nu-d\mid x]=-d,
    \qquad
    \mathbb E[(\nu-d)^2\mid x]=d^2+\sigma_D^2.
\end{equation}
Taking expectations and subtracting the result from $R(e)$ yields
\eqref{eq:exact-adaptation}.
\end{proof}

Equation~\eqref{eq:exact-adaptation} is the central algebraic result of the
minimal model. The two quantities
$\Delta_{\mathrm{now}}$ and $\Delta_R$ describe different effects of the
same routing action. The former measures the expected quality gain of the
current operational response, whereas the latter measures the expected
change in the future population risk of the cheap predictor induced by
training on that response.

Since
\begin{equation}
    \nabla R(e)=2Me,
    \qquad
    x^\top Me=\frac{1}{2}x^\top\nabla R(e),
\end{equation}
the first-order adaptation effect for small $\eta$ is governed by
\begin{equation}
    (\alpha-\beta)x^\top Me.
    \label{eq:geometric-factor}
\end{equation}
Thus, local teacher quality and training value are not equivalent: a teacher
response can improve the current prediction while inducing an update that
increases population risk, or conversely.

\subsection{Compatibility under isotropic risk geometry}

The first result identifies a setting in which immediate superiority and
positive adaptation value cannot conflict under a conservative update.

\begin{adaptationproposition}[Isotropic compatibility]
Suppose $M=cI$ with $c>0$ and define
\begin{equation}
    a:=\eta\|x\|^2.
\end{equation}
If
\begin{equation}
    0<a\leq1
    \qquad\text{and}\qquad
    \Delta_{\mathrm{now}}>0,
\end{equation}
then
\begin{equation}
    \Delta_R>0.
\end{equation}
\end{adaptationproposition}

\begin{proof}
From $\Delta_{\mathrm{now}}>0$,
\begin{equation}
    \sigma_D^2<\alpha^2-\beta^2.
    \label{eq:teacher-better-condition}
\end{equation}
In particular, $|\alpha|>|\beta|$, which implies
\begin{equation}
    \alpha(\alpha-\beta)>0.
\end{equation}
Under $M=cI$, equation~\eqref{eq:exact-adaptation} becomes
\begin{equation}
    \Delta_R
    =
    c\eta
    \left[
        2\alpha(\alpha-\beta)
        -a\bigl\{(\alpha-\beta)^2+\sigma_D^2\bigr\}
    \right].
\end{equation}
Using~\eqref{eq:teacher-better-condition},
\begin{align}
    (\alpha-\beta)^2+\sigma_D^2
    &<
    (\alpha-\beta)^2+\alpha^2-\beta^2
    \notag\\
    &=
    2\alpha(\alpha-\beta).
\end{align}
Since $a>0$,
\begin{align}
    \Delta_R
    &>
    2c\eta\alpha(\alpha-\beta)(1-a).
    \label{eq:isotropic-bound}
\end{align}
For $0<a<1$, the right-hand side is strictly positive because
$\alpha(\alpha-\beta)>0$. For $a=1$, the strict inequality in
\eqref{eq:isotropic-bound} directly gives $\Delta_R>0$.
\end{proof}

The condition $0<a\leq1$ will be referred to as the conservative
no-overshoot regime. For a repeated fixed input, the noiseless residual
multiplier is $1-a$, whose magnitude is below one for $0<a<2$.
The interval $1<a<2$ is therefore locally stable in this restricted sense,
but permits overshoot. Such local stability does not imply a decrease in
population risk in the presence of noise, anisotropic risk geometry, or
changing inputs. The compatibility result above is therefore not extended
to the entire interval $0<a<2$.

\subsection{Geometric conflict under anisotropic risk}

The isotropic result does not extend universally to anisotropic
second-moment matrices.

\begin{theorem}[Anisotropic conflict; preliminary]
Let $M\succ0$ be symmetric and not proportional to the identity. Then there
exist $x$ and $e$, together with a noisy teacher unbiased with respect to
the conditional mean
\begin{equation}
    b=0,
    \qquad
    0<\sigma_D^2<\alpha^2,
\end{equation}
such that
\begin{equation}
    \Delta_{\mathrm{now}}>0
    \qquad\text{and}\qquad
    \Delta_R<0
\end{equation}
for every $\eta>0$.
\end{theorem}

\begin{proof}
Because $M$ is not proportional to the identity, there exists a nonzero
$x$ that is not an eigenvector of $M$. Hence $x$ and $Mx$ are not
collinear, and strict Cauchy--Schwarz gives
\begin{equation}
    (x^\top Mx)^2
    <
    \|x\|^2\|Mx\|^2.
\end{equation}
Since $x^\top Mx>0$, the interval
\begin{equation}
    \frac{x^\top Mx}{\|Mx\|^2}
    <
    k
    <
    \frac{\|x\|^2}{x^\top Mx}
    \label{eq:k-interval}
\end{equation}
is nonempty.

Choose any $k$ satisfying~\eqref{eq:k-interval} and set
\begin{equation}
    e=x-kMx.
\end{equation}
Then
\begin{align}
    \alpha=x^\top e
    &=\|x\|^2-kx^\top Mx>0,
    \\
    x^\top Me
    &=x^\top Mx-k\|Mx\|^2<0.
\end{align}
Now take $b=0$, so that $\beta=0$, and choose
\begin{equation}
    0<\sigma_D^2<\alpha^2.
\end{equation}
Equation~\eqref{eq:immediate-gain} then gives
\begin{equation}
    \Delta_{\mathrm{now}}
    =\alpha^2-\sigma_D^2>0.
\end{equation}
On the other hand,
\begin{equation}
    \Delta_R
    =
    2\eta\alpha x^\top Me
    -
    \eta^2(\alpha^2+\sigma_D^2)x^\top Mx.
\end{equation}
The first term is strictly negative because
$\alpha>0$ and $x^\top Me<0$, while the second is strictly negative because
$M\succ0$ and $\eta>0$. Therefore
\begin{equation}
    \Delta_R<0
\end{equation}
for every $\eta>0$.
\end{proof}

The theorem is an existence result. It does not state that every query under
anisotropic risk geometry exhibits a conflict. Rather, isotropy guarantees
compatibility in the conservative regime, whereas anisotropy permits
configurations in which a teacher that improves the current expected
prediction induces a population-risk-increasing update, even for
arbitrarily small step sizes.

The construction is algebraic and does not imply that the conflicting
inputs have positive probability under every distribution sharing the same
second-moment matrix $M$. Establishing distributional conditions under which
such conflict occurs with nonzero probability is outside the present
single-update analysis.

Neither harmful updates nor gradient misalignment by themselves are claimed
as novel contributions. Their relevance here is that they can alter the
value of an operational routing decision when the same costly response is
also used to adapt the cheap predictor.

\subsection{Isolated-horizon routing comparison}

A purely myopic routing rule queries the costly predictor whenever
\begin{equation}
    \Delta_{\mathrm{now}}>C_D.
    \label{eq:myopic-rule}
\end{equation}

To expose the possible consequence of adaptation without yet solving the
full delayed-feedback problem, consider an idealized isolated horizon.
Suppose that the only future difference caused by the current routing
decision is a constant population-risk gap $\Delta_R$ over the next $H$
pre-feedback responses. Assume, in particular, no intervening updates and
no action-dependent differences in subsequent routing costs.

With discount factor $\gamma$, define
\begin{equation}
    B_H=\sum_{k=1}^{H}\gamma^k.
\end{equation}
Under these assumptions, the accumulated adaptation value is
\begin{equation}
    \Delta_{\mathrm{adapt}}^{(H)}
    =B_H\Delta_R,
\end{equation}
and the corresponding future-aware comparison becomes
\begin{equation}
    \Delta_{\mathrm{now}}+B_H\Delta_R>C_D.
    \label{eq:isolated-rule}
\end{equation}

This criterion is an analytical device rather than a solution of the
delayed-feedback objective. In particular, it does not assume that the
first reliable correction necessarily erases the difference between the
two routing histories, nor that all effects disappear after feedback
arrives.

\paragraph{Isolated over-querying.}
Consider a configuration satisfying
\begin{equation}
    \Delta_{\mathrm{now}}>0,
    \qquad
    \Delta_R<0,
    \qquad
    B_H>0.
\end{equation}
For any nonnegative cost satisfying
\begin{equation}
    \max\{0,\Delta_{\mathrm{now}}+B_H\Delta_R\}
    <
    C_D
    <
    \Delta_{\mathrm{now}},
    \label{eq:over-querying-cost}
\end{equation}
the myopic rule~\eqref{eq:myopic-rule} queries the teacher whereas the
isolated-horizon rule~\eqref{eq:isolated-rule} does not. The interval in
\eqref{eq:over-querying-cost} is nonempty because
$\Delta_{\mathrm{now}}>0$ and $\Delta_R<0$.

This is an existence result for suitable routing costs under the isolated
comparison, not a statement about every cost or about the optimal sequential
policy.

Conversely, isolated under-querying occurs whenever
\begin{equation}
    \Delta_{\mathrm{now}}
    <
    C_D
    <
    \Delta_{\mathrm{now}}+B_H\Delta_R.
\end{equation}
This requires $\Delta_R>0$ and can already occur in the scalar model.
Under-querying and over-querying are therefore distinct mechanisms and are
not asserted here as a single general sequential theorem.

\subsection{Open problem: counterfactual evolution during delayed feedback}

The isolated-horizon approximation suppresses the central sequential
difficulty: after the initial routing decision, the two systems may undergo
different intermediate updates and routing decisions before the reliable
target $Y_t$ becomes available.

Starting from the same pre-query state, let $e_{t+k}^{F}$ and
$e_{t+k}^{D}$ denote the cheap-predictor errors under the two
counterfactual histories generated by the initial choices $F$ and $D$,
respectively. Intervening updates and delayed reliable feedback from earlier
samples are allowed.

A natural population-risk component of the accumulated adaptation value is
\begin{equation}
    \Delta_{\mathrm{adapt}}^{(\tau)}
    =
    \sum_{k=1}^{\tau}
    \gamma^{k-1}
    \mathbb E\!\left[
        (e_{t+k}^{F})^\top M e_{t+k}^{F}
        -
        (e_{t+k}^{D})^\top M e_{t+k}^{D}
    \right].
    \label{eq:counterfactual-adaptation}
\end{equation}
The discount origin in~\eqref{eq:counterfactual-adaptation} is the first
future response. Hence its contribution measured from the current round is
\begin{equation}
    \gamma\Delta_{\mathrm{adapt}}^{(\tau)}.
\end{equation}
Under the isolated assumptions with $H=\tau$, this reduces to
\begin{equation}
    \gamma\Delta_{\mathrm{adapt}}^{(\tau)}
    =
    B_H\Delta_R.
\end{equation}
For $\gamma=0$, future adaptation has zero weight in the current objective.

Equation~\eqref{eq:counterfactual-adaptation} does not yet solve the
sequential routing problem. The correction dynamics, intermediate update
rules, and continuation policies must be specified before the two
counterfactual trajectories can be evaluated. Moreover, the expression
captures only differences in the cheap predictor's population risk; the
complete continuation-value difference must also include future operational
responses and consultation costs.

No regret bound, closed-form solution of the counterfactual trajectories,
or optimality claim for the full delayed-feedback policy is made at this
stage.